% ======================================================================
% diagrams.tex — Sequence Diagrams and Class Diagrams for ExamGen
% ======================================================================

\newpage
\subsection*{UML Diagrams}
\addcontentsline{toc}{subsection}{UML Diagrams}
\rhead{\itshape UML Diagrams}

% ======================================================================
% SEQUENCE DIAGRAM 1 — PDF Upload and docParser Pipeline
% ======================================================================

\subsubsection*{SD1: PDF Upload and Parsing Pipeline}

\begin{figure}[H]
\centering
\resizebox{0.88\textwidth}{!}{%
\begin{sequencediagram}

\newthread[white]{t}{\textbf{Teacher}}
\newinst[2.2]{web}{\textbf{WebApp}}
\newinst[2.2]{dp}{\textbf{docParser}}
\newinst[2.2]{ef}{\textbf{exam-forge}}

\begin{call}{t}{uploadPDF(file)}{web}{status: parsed}
  \begin{call}{web}{POST /courses/upload}{web}{}
  \end{call}
  \begin{call}{web}{parseAsync(file\_path)}{dp}{JSON saved}
    \begin{call}{dp}{\textbf{S1}: Docling $\to$ Markdown + PyMuPDF fill}{dp}{}
    \end{call}
    \begin{call}{dp}{\textbf{S2}: LLM + Instructor $\to$ Exam JSON}{dp}{}
    \end{call}
    \begin{call}{dp}{\textbf{S3}: Grounding validator + confidence $c$}{dp}{}
    \end{call}
    \begin{call}{dp}{\textbf{S4}: GPT-OSS-120B $\to$ Bloom level}{dp}{}
    \end{call}
    \begin{call}{dp}{\textbf{S5}: Qwen2.5-VL $\to$ diagram transcription}{dp}{}
    \end{call}
  \end{call}
  \begin{call}{web}{POST /api/invalidate-cache}{ef}{ok}
  \end{call}
  \begin{call}{web}{POST /api/index-file}{ef}{ok}
  \end{call}
\end{call}

\end{sequencediagram}
}
\caption{SD1: PDF upload and \texttt{docParser} five-stage IDP pipeline.}
\label{fig:seq_upload}
\end{figure}

% ======================================================================
% SEQUENCE DIAGRAM 2 — Exam Generation with RAG and Evaluation
% ======================================================================

\subsubsection*{SD2: Exam Generation with RAG and Post-Generation Evaluation}

\begin{figure}[H]
\centering
\resizebox{0.88\textwidth}{!}{%
\begin{sequencediagram}

\newthread[white]{t}{\textbf{Teacher}}
\newinst[1.6]{fe}{\textbf{Frontend}}
\newinst[1.6]{be}{\textbf{Backend}}
\newinst[1.6]{rag}{\textbf{RAGEngine}}
\newinst[1.6]{gen}{\textbf{Generator}}
\newinst[1.6]{ev}{\textbf{Evaluator}}

\begin{call}{t}{generateExam(subj, type, level)}{fe}{SSE stream}
  \begin{call}{fe}{POST /api/generate}{be}{}
    \begin{call}{be}{retrieve(query, subject)}{rag}{top-K chunks}
    \end{call}
    \begin{call}{be}{get\_style\_examples(subj, level)}{be}{few-shot anchors}
    \end{call}
    \begin{call}{be}{generate(context, style, type)}{gen}{question JSON}
    \end{call}
    \begin{call}{be}{evaluate(question)}{ev}{eval\_result}
      \begin{call}{ev}{fresh\_retrieve(question\_stem)}{rag}{fresh context}
      \end{call}
    \end{call}
    \mess{be}{SSE: question + evaluation}{fe}
  \end{call}
\end{call}
\mess{fe}{card: accept / edit / reject}{t}

\end{sequencediagram}
}
\caption{SD2: Exam generation with hybrid RAG retrieval and independent post-generation evaluation.}
\label{fig:seq_gen}
\end{figure}

% ======================================================================
% CLASS DIAGRAM 1 — docParser Data Schema
% ======================================================================

\subsubsection*{CD1: docParser Data Schema}

\begin{figure}[H]
\centering
\begin{tikzpicture}[
  every node/.style={font=\small},
  cls/.style={
    draw=black!80, thick,
    rectangle split,
    rectangle split parts=3,
    rectangle split part fill={blue!25, blue!5, white},
    rounded corners=2pt,
    text width=4.0cm,
    align=left
  },
  comp/.style={->, >=latex, thick, black},
  dep/.style={->, >=latex, thick, black!50, dashed}
]

% --- Exam ---
\node[cls] (exam) at (0,0) {
  \centering\textbf{Exam}
  \nodepart{two}
  \texttt{subject : str}\\
  \texttt{source\_file : str}\\
  \texttt{confidence : float}\\
  \texttt{needs\_review : bool}
  \nodepart{three}
  ~
};

% --- Exercise ---
\node[cls, below=1.4cm of exam] (exercise) {
  \centering\textbf{Exercise}
  \nodepart{two}
  \texttt{title : str}\\
  \texttt{context : str}\\
  \texttt{points : float}
  \nodepart{three}
  ~
};

% --- Question ---
\node[cls, below=1.4cm of exercise] (question) {
  \centering\textbf{Question}
  \nodepart{two}
  \texttt{question\_text : str}\\
  \texttt{type : QCM|FREE|CODE}\\
  \texttt{bloom\_level : str}\\
  \texttt{points : float}
  \nodepart{three}
  ~
};

% --- Solution ---
\node[cls, right=2.8cm of question] (solution) {
  \centering\textbf{Solution}
  \nodepart{two}
  \texttt{answer\_text : str}\\
  \texttt{choices : str[]}\\
  \texttt{correct\_choice : int}\\
  \texttt{grading\_rubric : str[]}
  \nodepart{three}
  ~
};

% --- TableData ---
\node[cls, above=0.7cm of solution] (tabledata) {
  \centering\textbf{TableData}
  \nodepart{two}
  \texttt{headers : str[]}\\
  \texttt{rows : str[][]}
  \nodepart{three}
  ~
};

% --- DiagramData ---
\node[cls, below=0.7cm of solution] (diagramdata) {
  \centering\textbf{DiagramData}
  \nodepart{two}
  \texttt{description : str}\\
  \texttt{diagram\_type : str}
  \nodepart{three}
  ~
};

% Composition arrows (left column)
\draw[comp] (exam.south) -- node[right, font=\scriptsize]{\texttt{1..*}} (exercise.north);
\draw[comp] (exercise.south) -- node[right, font=\scriptsize]{\texttt{1..*}} (question.north);

% Dependency arrows (right, optional fields)
\draw[dep] (question.east) -- node[above, font=\scriptsize]{\texttt{0..1}} (solution.west);
\draw[dep] (question.east) -- node[above, font=\scriptsize]{\texttt{0..1}} (tabledata.west);
\draw[dep] (question.east) -- node[below, font=\scriptsize]{\texttt{0..1}} (diagramdata.west);

\end{tikzpicture}
\caption{CD1: \texttt{docParser} data schema hierarchy
(\texttt{Exam} $\supset$ \texttt{Exercise} $\supset$ \texttt{Question}).}
\label{fig:class_schema}
\end{figure}

% ======================================================================
% CLASS DIAGRAM 2 — System Architecture Components
% ======================================================================

\subsubsection*{CD2: Main System Components}

\begin{figure}[H]
\centering
\begin{tikzpicture}[
  every node/.style={font=\small},
  cls/.style={
    draw=black!80, thick,
    rectangle split,
    rectangle split parts=3,
    rectangle split part fill={green!25, green!5, white},
    rounded corners=2pt,
    text width=4.6cm,
    align=left
  },
  uses/.style={->, >=latex, thick, dashed, gray!70}
]

% --- Dataset (top-left) ---
\node[cls] (dataset) at (0,0) {
  \centering\textbf{Dataset}
  \nodepart{two}
  \texttt{teacher\_id : str}\\
  \texttt{subject : str}\\
  \texttt{data\_path : str}
  \nodepart{three}
  \texttt{load\_questions()}\\
  \texttt{get\_examples(s, lvl)}\\
  \texttt{invalidate\_cache()}
};

% --- RAGEngine (top-right) ---
\node[cls, right=2.2cm of dataset] (rag) {
  \centering\textbf{RAGEngine}
  \nodepart{two}
  \texttt{bm25 : BM25Index}\\
  \texttt{faiss\_ar : FAISSIndex}\\
  \texttt{faiss\_ml : FAISSIndex}
  \nodepart{three}
  \texttt{retrieve(q, subj, k)}\\
  \texttt{route\_language(q)}\\
  \texttt{index\_file(path)}
};

% --- QuestionGenerator (bottom-left) ---
\node[cls, below=2.0cm of dataset] (gen) {
  \centering\textbf{QuestionGenerator}
  \nodepart{two}
  \texttt{model : str}\\
  \texttt{temperature : float}\\
  \texttt{max\_retries : int}
  \nodepart{three}
  \texttt{generate\_mcq(ctx, style)}\\
  \texttt{generate\_saq(ctx, style)}\\
  \texttt{generate\_exercise(ctx)}
};

% --- PostEvaluator (bottom-right) ---
\node[cls, right=2.2cm of gen] (ev) {
  \centering\textbf{PostEvaluator}
  \nodepart{two}
  \texttt{rag : RAGEngine}\\
  \texttt{judge\_model : str}
  \nodepart{three}
  \texttt{evaluate(question)}\\
  \texttt{fresh\_retrieve(stem)}\\
  ~
};

% Arrows
\draw[uses] (gen.north) -- node[left, font=\scriptsize]{uses} (dataset.south);
\draw[uses] (gen.east) -- node[above, font=\scriptsize]{uses} (ev.west);
\draw[uses] (ev.north) -- node[right, font=\scriptsize]{uses} (rag.south);
\draw[uses] ([xshift=1cm]gen.north) -- ([xshift=1cm]dataset.south -| rag.south)
  node[midway, right, font=\scriptsize]{uses} (rag.south);

\end{tikzpicture}
\caption{CD2: Main system components and their dependencies.}
\label{fig:class_system}
\end{figure}
